\section{\texorpdfstring{$\zeta$}{zeta} 型对象与有限部：离散读出的解析接口}

\subsection{动力学 \texorpdfstring{$\zeta$}{zeta} 函数（稳定语法层）}

在稳定语法层，考虑黄金均值约束的双边移位空间 $X\subset\{0,1\}^{\ZZ}$（禁止相邻 $11$）及其左移位 $\sigma$。对一般动力系统 $(Y,S)$，动力学 $\zeta$ 函数定义为
\[
\zeta_{S}(z):=\exp\left(\sum_{n\ge1}\frac{\#\mathrm{Fix}(S^n)}{n}z^n\right),
\]
其中 $\mathrm{Fix}(S^n)$ 为 $n$ 周期点集合。对有限型子移位（由邻接矩阵 $A$ 给出），有标准恒等式
\[
\zeta_{\sigma}(z)=\frac{1}{\det(I-zA)}
\quad\text{且}\quad
\#\mathrm{Fix}(\sigma^n)=\mathrm{Tr}(A^n)
\]
（参见 \cite{LindMarcus1995} 中关于有限型子移位与 $\zeta$ 函数的标准章节）。对黄金均值约束的
\(
A=\begin{pmatrix}1&1\\1&0\end{pmatrix}
\)
即得到
\[
\zeta(z)=\frac{1}{\det(I-zA)}=\frac{1}{1-z-z^2}.
\]
其极点结构直接由 $\varphi$ 控制。这提供一个“语法—谱—$\zeta$”的闭合环：黄金比例既决定语法的增长率，也决定 $\zeta$ 的主极点。

\subsection{轨道 Dirichlet 级数与 Abel 有限部（扫描层）}

在扫描—投影框架中，另一个自然 $\zeta$ 型对象来自对轨道读出序列的加权求和。设 $f:X\to\CC$ 为可测函数（例如 $f=\ind_W$ 或其中心化版本），定义
\[
L(f,s;x)=\sum_{n\ge1}\frac{f(T^n x)}{n^s}.
\]
当 $f\equiv 1$ 时退化为经典 $\zeta(s)$。对一般 $f$，可通过 Abel 求和与解析延拓讨论其收敛域外的有限部。本文强调的结构性观点是：
\begin{itemize}
  \item 离散读出与有限分辨率折叠导致大量求和对象在临界点附近发散；
  \item 若将“读出”理解为一种投影仪器，则有限部提取可被视为与折叠相配套的解析层读出规则；
  \item 因而“离散化—正则化—有限部”可作为生成链中与算术编码并列的稳定接口。
\end{itemize}
\begin{definition}[Abel 正则化与 Hadamard 有限部]\label{def:abel-finite-part}
对任意 $s\in\CC$，若极限存在，定义轨道 Dirichlet 级数的 Abel 正则化为
\[
L_{\mathrm{Abel}}(f,s;x):=\lim_{r\uparrow 1}\ \sum_{n\ge1}\frac{f(T^n x)\,r^n}{n^s}.
\]
若 $L_{\mathrm{Abel}}(f,s;x)$ 在点 $s_0$ 的某邻域内可作 Laurent 展开
\[
L_{\mathrm{Abel}}(f,s;x)=\sum_{k=-K}^{\infty} a_k\,(s-s_0)^k,
\]
则称常数项 $a_0$ 为其在 $s_0$ 的 Hadamard 有限部，记为 $\mathrm{fp}_{s=s_0}L_{\mathrm{Abel}}(f,s;x):=a_0$。
\end{definition}

定义 \ref{def:abel-finite-part} 所用仅为 Abel 求和与解析延拓的标准工具，见 \cite{Apostol1976}。

\begin{proposition}[最小存在性口径]\label{prop:abel-finite-part-minimal}
若 $f$ 有界，则对任意 $r\in(0,1)$ 级数
\(
\sum_{n\ge 1} f(T^n x)\,r^n/n^s
\)
在 $\Re(s)>0$ 绝对收敛并定义解析函数；当 $\Re(s)>1$ 时，$r\uparrow 1$ 的极限（若存在）与通常的 Dirichlet 级数 $L(f,s;x)$ 一致。若进一步已知 $L_{\mathrm{Abel}}(f,s;x)$ 在某点 $s_0$ 可作 Laurent 展开，则有限部 $\mathrm{fp}_{s=s_0}$ 作为常数项是标准且不依赖展开方式的对象（见 \cite{Apostol1976}）。
\end{proposition}

